\documentclass[12pt,a4paper]{article}
\usepackage[UTF8]{ctex}
\usepackage{geometry}
\usepackage{listings}
\usepackage{xcolor}
\usepackage{hyperref}
\usepackage{booktabs}
\usepackage{float}
\usepackage{tcolorbox}
\usepackage{enumitem}

\geometry{left=2.5cm,right=2.5cm,top=2.5cm,bottom=2.5cm}

\lstset{
    language=Python,
    basicstyle=\ttfamily\small,
    keywordstyle=\color{blue}\bfseries,
    commentstyle=\color{gray},
    stringstyle=\color{orange},
    numbers=left,
    numberstyle=\tiny\color{gray},
    stepnumber=1,
    numbersep=8pt,
    backgroundcolor=\color{gray!10},
    frame=single,
    frameround=fttt,
    breaklines=true,
    showstringspaces=false,
    tabsize=4,
    captionpos=b,
    xleftmargin=2em,
    xrightmargin=1em,
    aboveskip=1em,
    belowskip=1em
}

\tcbuselibrary{skins,breakable}
\newtcolorbox{knowledgebox}[1][]{
    colback=blue!5!white,
    colframe=blue!50!black,
    fonttitle=\bfseries,
    title=#1,
    breakable
}

\newtcolorbox{practicebox}[1][]{
    colback=green!5!white,
    colframe=green!50!black,
    fonttitle=\bfseries,
    title=#1,
    breakable
}

\newtcolorbox{tipbox}[1][]{
    colback=yellow!10!white,
    colframe=orange!70!black,
    fonttitle=\bfseries,
    title=#1,
    breakable
}

\hypersetup{
    colorlinks=true,
    linkcolor=blue,
    filecolor=magenta,      
    urlcolor=cyan,
    bookmarks=true,
    pdfstartview=FitH
}

\title{\textbf{Pandas零基础入门手册}\\[0.5em]\large 从Python基础到数据处理实战}
\author{竞赛培训组}
\date{\today}

\begin{document}

\maketitle
\tableofcontents
\newpage

%========================================
\section{Python基础入门}
%========================================

\subsection{什么是编程？}

\begin{knowledgebox}[知识点：编程是什么？]
编程就是给计算机写指令，让计算机按照我们的要求去工作。

\textbf{类比理解：}
\begin{itemize}
    \item 编程就像写菜谱：步骤清晰、顺序明确
    \item 计算机就像厨师：严格按照菜谱操作
    \item 程序员就像菜谱作者：设计制作流程
\end{itemize}

\textbf{为什么要学编程？}
\begin{itemize}
    \item 自动化处理重复工作
    \item 处理大量数据（Excel处理不了的）
    \item 竞赛必备技能
\end{itemize}
\end{knowledgebox}

\subsection{Python是什么？}

Python是一种\textbf{编程语言}，就像中文、英语是人类交流的语言一样，Python是人和计算机交流的语言。

\begin{knowledgebox}[Python的特点]
\begin{enumerate}
    \item \textbf{简单易学}
    \begin{itemize}
        \item 语法接近自然语言
        \item 代码可读性强
        \item 适合初学者入门
    \end{itemize}
    
    \item \textbf{功能强大}
    \begin{itemize}
        \item 数据分析（Pandas）
        \item 人工智能（TensorFlow）
        \item 网站开发（Django）
        \item 自动化办公
    \end{itemize}
    
    \item \textbf{应用广泛}
    \begin{itemize}
        \item 大数据处理
        \item 机器学习
        \item 科学研究
        \item 金融分析
    \end{itemize}
\end{enumerate}
\end{knowledgebox}

\subsection{安装Python}

\textbf{步骤：}
\begin{enumerate}
    \item 访问官网：https://www.python.org/downloads/
    \item 下载Python 3.8或更高版本
    \item 运行安装程序，\textbf{务必勾选"Add Python to PATH"}
    \item 点击Install Now完成安装
\end{enumerate}

\begin{tipbox}[重要提示]
安装时一定要勾选"Add Python to PATH"！

如果不勾选，后续会出现"python不是内部或外部命令"的错误。
\end{tipbox}

\subsection{验证安装}

\begin{enumerate}
    \item 按Win+R，输入cmd，回车
    \item 输入：\texttt{python --version}
    \item 如果显示版本号（如Python 3.9.0），说明安装成功
\end{enumerate}

%========================================
\section{Python基础语法详解}
%========================================

\subsection{变量和数据类型}

\begin{knowledgebox}[什么是变量？]
变量就是给数据起个名字，方便以后使用。

\textbf{类比理解：}
\begin{itemize}
    \item 变量名 = 标签（如"年龄"）
    \item 变量值 = 实际数据（如20）
    \item 赋值 = 把数据存到这个名字里
\end{itemize}

\textbf{示例：}
\begin{lstlisting}
age = 20  # 把数字20存到名字"age"里
name = "张三"  # 把文字"张三"存到名字"name"里
\end{lstlisting}
\end{knowledgebox}

\textbf{Python常用数据类型：}

\begin{table}[H]
\centering
\caption{Python数据类型}
\begin{tabular}{|l|l|l|l|}
\hline
\textbf{类型} & \textbf{名称} & \textbf{说明} & \textbf{示例} \\
\hline
int & 整数 & 没有小数的数字 & 20, 100, -5 \\
\hline
float & 浮点数 & 带小数的数字 & 1.75, 3.14 \\
\hline
str & 字符串 & 文字或字符 & "张三", "hello" \\
\hline
bool & 布尔值 & 真或假 & True, False \\
\hline
list & 列表 & 一组数据 & [1, 2, 3] \\
\hline
dict & 字典 & 键值对数据 & \{"name": "张三"\} \\
\hline
\end{tabular}
\end{table}

\subsection{列表（List）详解}

\begin{knowledgebox}[什么是列表？]
列表是一组数据的集合，用方括号[]表示。

\textbf{类比理解：}
\begin{itemize}
    \item 列表 = 购物清单、成绩单、通讯录
    \item 可以存放多个数据
    \item 数据之间有顺序（第1个、第2个...）
\end{itemize}
\end{knowledgebox}

\begin{lstlisting}[caption={列表示例}]
# 创建一个列表
scores = [85, 92, 78, 90, 88]
names = ["张三", "李四", "王五"]

# 访问列表元素（索引从0开始）
print(scores[0])    # 第一个元素：85
print(scores[1])    # 第二个元素：92
print(names[2])     # 第三个元素：王五

# 获取列表长度（元素个数）
print(len(scores))  # 5

# 添加元素到列表末尾
scores.append(95)
print(scores)       # [85, 92, 78, 90, 88, 95]
\end{lstlisting}

\begin{tipbox}[索引从0开始]
Python中，列表的索引从0开始：
\begin{itemize}
    \item 第1个元素的索引是0
    \item 第2个元素的索引是1
    \item 第n个元素的索引是n-1
\end{itemize}

这是编程界的惯例，几乎所有编程语言都是这样。
\end{tipbox}

\subsection{字典（Dictionary）详解}

\begin{knowledgebox}[什么是字典？]
字典是一种键值对的数据结构，用花括号\{\}表示。

\textbf{类比理解：}
\begin{itemize}
    \item 字典 = 真实字典（查词找释义）
    \item 键（key）= 要查的词（如"姓名"）
    \item 值（value）= 词的释义（如"张三"）
\end{itemize}
\end{knowledgebox}

\begin{lstlisting}[caption={字典示例}]
# 创建一个字典
student = {
    "name": "张三",
    "age": 20,
    "score": 85
}

# 访问字典元素（通过键访问值）
print(student["name"])   # 张三
print(student["age"])    # 20

# 修改值
student["score"] = 90
print(student["score"])  # 90

# 添加新键值对
student["gender"] = "男"
print(student)  # {"name": "张三", "age": 20, "score": 90, "gender": "男"}
\end{lstlisting}

\subsection{条件判断（if语句）}

\begin{knowledgebox}[什么是条件判断？]
条件判断就是根据不同的情况执行不同的代码。

\textbf{类比理解：}
\begin{itemize}
    \item 如果下雨，就带伞
    \item 如果成绩>=60，就及格
    \item 如果年龄>=18，就是成年人
\end{itemize}
\end{knowledgebox}

\begin{lstlisting}[caption={条件判断示例}]
# 基础if-else
age = 18

if age >= 18:
    print("成年人")
else:
    print("未成年人")

# 多条件判断（if-elif-else）
score = 85

if score >= 90:
    print("优秀")
elif score >= 80:
    print("良好")
elif score >= 60:
    print("及格")
else:
    print("不及格")
\end{lstlisting}

\begin{tipbox}[注意缩进]
Python用\textbf{缩进}（空格或Tab）来表示代码块，不是大括号！

\textbf{正确：}
\begin{lstlisting}
if age >= 18:
    print("成年人")  # 缩进4个空格
    print("可以考驾照")
\end{lstlisting}

\textbf{错误：}
\begin{lstlisting}
if age >= 18:
print("成年人")  # 没有缩进，会报错！
\end{lstlisting}
\end{tipbox}

\subsection{循环（for循环）}

\begin{knowledgebox}[什么是循环？]
循环就是重复执行某段代码。

\textbf{类比理解：}
\begin{itemize}
    \item 循环 = 跑步机（重复跑步动作）
    \item 循环 = 洗衣机（重复洗涤动作）
    \item 循环 = 打卡（每天重复打卡）
\end{itemize}
\end{knowledgebox}

\begin{lstlisting}[caption={循环示例}]
# 遍历列表
scores = [85, 92, 78, 90]

for score in scores:
    print(score)  # 依次打印每个分数

# 遍历范围（0到4）
for i in range(5):
    print(i)  # 打印0, 1, 2, 3, 4

# 遍历字典
student = {"name": "张三", "age": 20}

for key, value in student.items():
    print(f"{key}: {value}")
    # 输出：
    # name: 张三
    # age: 20
\end{lstlisting}

\subsection{函数（Function）}

\begin{knowledgebox}[什么是函数？]
函数是一段可以重复使用的代码块。

\textbf{类比理解：}
\begin{itemize}
    \item 函数 = 洗衣机、电饭煲（封装好的功能）
    \item 输入 = 脏衣服、米（参数）
    \item 输出 = 干净衣服、米饭（返回值）
    \item 好处 = 避免重复写代码
\end{itemize}
\end{knowledgebox}

\begin{lstlisting}[caption={函数示例}]
# 定义函数（无参数）
def say_hello():
    """打招呼函数"""
    print("你好！")

# 调用函数
say_hello()  # 输出：你好！

# 定义函数（有参数）
def greet(name):
    """带参数的打招呼函数"""
    print(f"你好，{name}！")

# 调用函数
 greet("张三")  # 输出：你好，张三！

# 定义函数（有返回值）
def add(a, b):
    """计算两数之和"""
    return a + b

# 调用函数
result = add(3, 5)
print(result)  # 输出：8
\end{lstlisting}

%========================================
\section{Pandas安装与基础}
%========================================

\subsection{什么是Pandas？}

\begin{knowledgebox}[Pandas是什么？]
Pandas是Python的一个数据分析库，专门用来处理表格数据。

\textbf{类比理解：}
\begin{itemize}
    \item Pandas = Excel的Python版本
    \item 但比Excel更强大、更灵活
    \item 可以处理几百万行数据（Excel会卡死）
\end{itemize}

\textbf{为什么叫Pandas？}
\begin{itemize}
    \item 和熊猫（panda）没有关系
    \item 来自"Panel Data"（面板数据）
    \item 是计量经济学中的术语
\end{itemize}
\end{knowledgebox}

\subsection{安装Pandas}

\begin{lstlisting}[caption={安装Pandas}]
# 在命令行中运行
pip install pandas

# 如果安装慢，使用国内镜像
pip install pandas -i https://pypi.tuna.tsinghua.edu.cn/simple
\end{lstlisting}

\subsection{导入Pandas}

\begin{lstlisting}[caption={导入Pandas}]
# 导入pandas，并起别名pd（行业惯例）
import pandas as pd

# 通常也导入numpy（数值计算库）
import numpy as np
\end{lstlisting}

\begin{tipbox}[为什么要用别名pd？]
\begin{itemize}
    \item \textbf{简化代码}：写pd比写pandas更短
    \item \textbf{行业惯例}：大家都用pd，方便阅读
    \item \textbf{保持一致}：不同人写的代码风格统一
\end{itemize}
\end{tipbox}

%========================================
\section{四个练习数据文件}
%========================================

\subsection{数据文件概览}

本手册所有示例代码都围绕以下\textbf{四个CSV文件}展开：

\begin{table}[H]
\centering
\caption{练习数据文件说明}
\begin{tabular}{|l|l|l|l|}
\hline
\textbf{文件名} & \textbf{数据量} & \textbf{主题} & \textbf{主要问题} \\
\hline
01\_电商订单数据\_练习.csv & 50行 & 电商订单 & 重复、缺失、异常值、格式混乱 \\
\hline
02\_传感器数据\_练习.csv & 144行 & 传感器读数 & 时间序列、格式问题 \\
\hline
03\_缺失值处理\_练习.csv & 30行 & 员工信息 & 各种缺失值场景 \\
\hline
04\_数据格式清洗\_练习.csv & 30行 & 用户信息 & 格式不统一、单位混乱 \\
\hline
\end{tabular}
\end{table}

\subsection{01\_电商订单数据\_练习.csv}

\begin{knowledgebox}[数据字段说明]
\begin{itemize}
    \item \textbf{order\_id}：订单编号
    \item \textbf{user\_id}：用户ID（部分缺失）
    \item \textbf{product\_name}：商品名称
    \item \textbf{price}：价格（带"元"单位）
    \item \textbf{quantity}：数量（有负数异常值）
    \item \textbf{order\_date}：订单日期（多种格式）
    \item \textbf{address}：地址
    \item \textbf{phone}：手机号（带横线）
    \item \textbf{tags}：商品标签（JSON格式）
\end{itemize}
\end{knowledgebox}

\subsection{02\_传感器数据\_练习.csv}

\begin{knowledgebox}[数据字段说明]
\begin{itemize}
    \item \textbf{timestamp}：时间戳
    \item \textbf{sensor\_id}：传感器ID
    \item \textbf{temperature}：温度
    \item \textbf{humidity}：湿度（部分缺失）
    \item \textbf{pressure}：压力
    \item \textbf{status}：设备状态
\end{itemize}
\end{knowledgebox}

\subsection{03\_缺失值处理\_练习.csv}

\begin{knowledgebox}[数据字段说明]
\begin{itemize}
    \item \textbf{id}：编号
    \item \textbf{name}：姓名
    \item \textbf{age}：年龄（部分缺失）
    \item \textbf{gender}：性别（部分缺失）
    \item \textbf{salary}：工资（部分缺失）
    \item \textbf{department}：部门（部分缺失）
    \item \textbf{join\_date}：入职日期（部分缺失）
    \item \textbf{phone}：电话
    \item \textbf{email}：邮箱
    \item \textbf{score}：评分（部分缺失）
\end{itemize}
\end{knowledgebox}

\subsection{04\_数据格式清洗\_练习.csv}

\begin{knowledgebox}[数据字段说明]
\begin{itemize}
    \item \textbf{id}：编号
    \item \textbf{name}：姓名（带空格）
    \item \textbf{birth\_date}：出生日期（多种格式）
    \item \textbf{height}：身高（带单位）
    \item \textbf{weight}：体重（带单位）
    \item \textbf{income}：收入（带k单位）
    \item \textbf{phone}：电话
    \item \textbf{address}：地址
    \item \textbf{product\_tags}：商品标签
\end{itemize}
\end{knowledgebox}

%========================================
\section{读取CSV文件}
%========================================

\subsection{什么是CSV文件？}

\begin{knowledgebox}[CSV文件介绍]
CSV（Comma-Separated Values，逗号分隔值）是一种简单的数据存储格式。

\textbf{特点：}
\begin{itemize}
    \item 纯文本格式，可以用记事本打开
    \item 每行是一条记录
    \item 每个字段用逗号分隔
    \item 第一行通常是列名（表头）
\end{itemize}

\textbf{示例：}
\begin{verbatim}
姓名,年龄,成绩
张三,20,85
李四,21,92
\end{verbatim}
\end{knowledgebox}

\subsection{读取四个练习文件}

\begin{lstlisting}[caption={读取CSV文件}]
import pandas as pd

# 读取电商订单数据
df_order = pd.read_csv('01_电商订单数据_练习.csv')

# 读取传感器数据
df_sensor = pd.read_csv('02_传感器数据_练习.csv')

# 读取缺失值练习数据
df_missing = pd.read_csv('03_缺失值处理_练习.csv')

# 读取格式清洗练习数据
df_messy = pd.read_csv('04_数据格式清洗_练习.csv')
\end{lstlisting}

\begin{knowledgebox}[知识点：pd.read\_csv()]
\texttt{pd.read\_csv()}是Pandas中最常用的函数，用于读取CSV文件。

\textbf{常用参数：}
\begin{itemize}
    \item \texttt{encoding='utf-8'}：指定文件编码（处理中文）
    \item \texttt{sep=','}：指定分隔符（默认是逗号）
    \item \texttt{header=0}：指定第几行是列名（默认第0行）
\end{itemize}
\end{knowledgebox}

\begin{practicebox}[动手练习1：读取并查看数据]
\textbf{任务：}读取四个CSV文件，查看每个文件的基本信息。

\begin{lstlisting}
import pandas as pd

# 读取四个文件
df_order = pd.read_csv('01_电商订单数据_练习.csv')
df_sensor = pd.read_csv('02_传感器数据_练习.csv')
df_missing = pd.read_csv('03_缺失值处理_练习.csv')
df_messy = pd.read_csv('04_数据格式清洗_练习.csv')

# 查看电商订单数据的形状
print("电商订单数据形状：", df_order.shape)
print("\n前5行：")
print(df_order.head())
\end{lstlisting}
\end{practicebox}

%========================================
\section{数据查看与探索}
%========================================

\subsection{查看数据基本信息}

\begin{lstlisting}[caption={查看数据形状和列名}]
# 查看数据形状（行数，列数）
print(df_order.shape)  # 输出：(50, 9)

# 查看列名
print(df_order.columns)

# 查看数据类型
print(df_order.dtypes)
\end{lstlisting}

\begin{knowledgebox}[知识点：shape属性]
\texttt{shape}返回一个元组\texttt{(行数, 列数)}，表示数据的大小。

\textbf{示例：}
\begin{itemize}
    \item \texttt{(50, 9)} 表示50行9列
    \item \texttt{df.shape[0]} 获取行数
    \item \texttt{df.shape[1]} 获取列数
\end{itemize}
\end{knowledgebox}

\subsection{查看数据内容}

\begin{lstlisting}[caption={查看数据内容}]
# 查看前5行
print(df_order.head())

# 查看后5行
print(df_order.tail())

# 查看前10行
print(df_order.head(10))
\end{lstlisting}

\subsection{统计信息}

\begin{lstlisting}[caption={使用describe()查看统计信息}]
# 查看数值列的统计信息
print(df_order.describe())

# 查看缺失值统计
print(df_order.isnull().sum())
\end{lstlisting}

\begin{knowledgebox}[知识点：describe()方法]
\texttt{describe()}方法返回数值列的统计摘要：
\begin{itemize}
    \item \texttt{count}：非空值数量
    \item \texttt{mean}：平均值
    \item \texttt{std}：标准差
    \item \texttt{min}：最小值
    \item \texttt{25\%}：第25百分位数（Q1）
    \item \texttt{50\%}：中位数
    \item \texttt{75\%}：第75百分位数（Q3）
    \item \texttt{max}：最大值
\end{itemize}
\end{knowledgebox}

\begin{practicebox}[动手练习2：探索传感器数据]
\textbf{任务：}使用02\_传感器数据\_练习.csv，查看数据的基本统计信息。

\begin{lstlisting}
import pandas as pd

df_sensor = pd.read_csv('02_传感器数据_练习.csv')

# 查看形状
print("数据形状：", df_sensor.shape)

# 查看统计信息
print("\n统计摘要：")
print(df_sensor.describe())

# 查看有多少个不同的传感器
print("\n传感器数量：", df_sensor['sensor_id'].nunique())
print("传感器列表：", df_sensor['sensor_id'].unique())
\end{lstlisting}
\end{practicebox}

%========================================
\section{数据选择}
%========================================

\subsection{选择列}

\begin{lstlisting}[caption={选择单列和多列}]
# 选择单列（商品名称）
products = df_order['product_name']

# 选择多列（订单ID、商品名称、价格）
order_info = df_order[['order_id', 'product_name', 'price']]
\end{lstlisting}

\begin{knowledgebox}[知识点：选择单列vs多列]
\begin{itemize}
    \item \textbf{单列}：\texttt{df['列名']}，返回Series
    \item \textbf{多列}：\texttt{df[['列1', '列2']]}，返回DataFrame
    \item \textbf{注意}：多列要用两个方括号！
\end{itemize}
\end{knowledgebox}

\subsection{选择行}

\begin{lstlisting}[caption={使用iloc选择行}]
# 选择第1行（索引为0）
row = df_order.iloc[0]

# 选择前5行
rows = df_order.iloc[0:5]

# 选择第10到20行
rows = df_order.iloc[10:20]
\end{lstlisting}

\begin{knowledgebox}[知识点：iloc vs loc]
\begin{itemize}
    \item \texttt{iloc}：按\textbf{位置}选择（用数字索引）
    \item \texttt{loc}：按\textbf{标签}选择（用行名）
    \item 新手建议先用\texttt{iloc}，更直观
\end{itemize}
\end{knowledgebox}

\subsection{条件筛选}

\begin{lstlisting}[caption={条件筛选示例}]
# 筛选数量大于1的订单
large_orders = df_order[df_order['quantity'] > 1]

# 筛选特定商品
iphone_orders = df_order[df_order['product_name'] == 'iPhone 14 Pro']
\end{lstlisting}

\begin{knowledgebox}[知识点：条件筛选原理]
条件筛选的原理：
\begin{enumerate}
    \item \texttt{df['列名'] > 值} 返回一个布尔Series（True/False）
    \item \texttt{df[布尔Series]} 返回True对应的行
\end{enumerate}

\textbf{常用条件运算符：}
\begin{itemize}
    \item \texttt{==}：等于
    \item \texttt{!=}：不等于
    \item \texttt{>}：大于
    \item \texttt{<}：小于
    \item \texttt{>=}：大于等于
    \item \texttt{<=}：小于等于
\end{itemize}
\end{knowledgebox}

\begin{practicebox}[动手练习3：筛选缺失值数据]
\textbf{任务：}使用03\_缺失值处理\_练习.csv，筛选出年龄缺失的员工。

\begin{lstlisting}
import pandas as pd

df_missing = pd.read_csv('03_缺失值处理_练习.csv')

# 筛选年龄缺失的行
missing_age = df_missing[df_missing['age'].isnull()]

print("年龄缺失的员工：")
print(missing_age[['name', 'age', 'department']])

print(f"\n共 {len(missing_age)} 人年龄缺失")
\end{lstlisting}
\end{practicebox}

%========================================
\section{数据清洗：处理缺失值}
%========================================

\subsection{什么是缺失值？}

\begin{knowledgebox}[缺失值介绍]
缺失值就是数据表中某些单元格没有数据，显示为NaN（Not a Number）或空值。

\textbf{产生原因：}
\begin{itemize}
    \item 数据采集时遗漏
    \item 用户没有填写
    \item 数据传输错误
    \item 某些字段不适用
\end{itemize}

\textbf{为什么要处理缺失值？}
\begin{itemize}
    \item 缺失值会影响统计计算（如平均值）
    \item 很多算法不接受缺失值
    \item 缺失值可能包含有用信息
\end{itemize}
\end{knowledgebox}

\subsection{检测缺失值}

\begin{lstlisting}[caption={检测缺失值}]
# 统计每列的缺失值数量
print(df_missing.isnull().sum())

# 计算缺失值比例
missing_ratio = df_missing.isnull().sum() / len(df_missing) * 100
print(missing_ratio)
\end{lstlisting}

\begin{knowledgebox}[知识点：isnull()方法]
\texttt{isnull()}方法返回一个布尔DataFrame：
\begin{itemize}
    \item \texttt{True}：表示这个单元格是缺失值
    \item \texttt{False}：表示这个单元格有数据
\end{itemize}

\texttt{sum()}会把True当成1、False当成0，所以能统计出每列的缺失值数量。
\end{knowledgebox}

\subsection{填充缺失值}

\begin{lstlisting}[caption={填充缺失值的方法}]
# 用固定值填充
df_missing['gender'] = df_missing['gender'].fillna('未知')

# 用平均值填充（工资列）
mean_salary = df_missing['salary'].mean()
df_missing['salary'] = df_missing['salary'].fillna(mean_salary)

# 用中位数填充（年龄列）
median_age = df_missing['age'].median()
df_missing['age'] = df_missing['age'].fillna(median_age)

# 用众数填充（部门列）
mode_dept = df_missing['department'].mode()[0]
df_missing['department'] = df_missing['department'].fillna(mode_dept)
\end{lstlisting}

\begin{knowledgebox}[知识点：填充策略选择]
\begin{itemize}
    \item \textbf{平均值}：适合正态分布的数值数据
    \item \textbf{中位数}：适合有异常值的数据（不受极端值影响）
    \item \textbf{众数}：适合分类数据（如性别、部门）
    \item \textbf{固定值}：如用"未知"填充
    \item \textbf{删除}：如果缺失值比例超过50\%，考虑删除该列
\end{itemize}
\end{knowledgebox}

\begin{practicebox}[动手练习4：处理缺失值]
\textbf{任务：}使用03\_缺失值处理\_练习.csv，完成缺失值处理。

\begin{lstlisting}
import pandas as pd

df = pd.read_csv('03_缺失值处理_练习.csv')

print("处理前缺失值统计：")
print(df.isnull().sum())

# 用中位数填充年龄
median_age = df['age'].median()
df['age'] = df['age'].fillna(median_age)

# 用众数填充性别
mode_gender = df['gender'].mode()[0]
df['gender'] = df['gender'].fillna(mode_gender)

# 用平均值填充工资
mean_salary = df['salary'].mean()
df['salary'] = df['salary'].fillna(mean_salary)

print("\n处理后缺失值统计：")
print(df.isnull().sum())
\end{lstlisting}
\end{practicebox}

%========================================
\section{数据清洗：处理重复值}
%========================================

\subsection{什么是重复值？}

\begin{knowledgebox}[重复值介绍]
重复值就是数据表中完全相同的记录出现了多次。

\textbf{产生原因：}
\begin{itemize}
    \item 数据录入重复
    \item 数据合并时产生
    \item 系统故障导致重复提交
\end{itemize}

\textbf{为什么要处理重复值？}
\begin{itemize}
    \item 重复值会导致统计结果不准确
    \item 浪费存储空间
    \item 影响数据分析结论
\end{itemize}
\end{knowledgebox}

\subsection{检测和删除重复值}

\begin{lstlisting}[caption={处理重复值}]
# 检测重复值
duplicates = df_order.duplicated()
print(f"重复行数：{duplicates.sum()}")

# 查看重复的行
duplicate_rows = df_order[df_order.duplicated(keep=False)]
print(duplicate_rows)

# 删除重复值（保留第一个）
df_order_clean = df_order.drop_duplicates()
\end{lstlisting}

\begin{knowledgebox}[知识点：duplicated()方法]
\texttt{duplicated()}方法返回一个布尔Series：
\begin{itemize}
    \item \texttt{keep='first'}：保留第一个，标记后面的为重复（默认）
    \item \texttt{keep='last'}：保留最后一个，标记前面的为重复
    \item \texttt{keep=False}：标记所有重复的行
\end{itemize}
\end{knowledgebox}

\begin{practicebox}[动手练习5：处理电商订单重复值]
\textbf{任务：}使用01\_电商订单数据\_练习.csv，检测并删除重复订单。

\begin{lstlisting}
import pandas as pd

df = pd.read_csv('01_电商订单数据_练习.csv')

print(f"原始数据行数：{len(df)}")

# 检测重复值
n_duplicates = df.duplicated().sum()
print(f"重复行数：{n_duplicates}")

# 查看重复的行
duplicate_rows = df[df.duplicated(keep=False)]
print("\n重复的行：")
print(duplicate_rows[['order_id', 'product_name', 'price']])

# 删除重复值
df_clean = df.drop_duplicates()
print(f"\n删除后行数：{len(df_clean)}")
\end{lstlisting}
\end{practicebox}

%========================================
\section{数据清洗：处理异常值}
%========================================

\subsection{什么是异常值？}

\begin{knowledgebox}[异常值介绍]
异常值是指明显偏离正常范围的数据点。

\textbf{示例：}
\begin{itemize}
    \item 年龄为负数或超过150岁
    \item 价格为0或异常高
    \item 数量为负数
\end{itemize}

\textbf{产生原因：}
\begin{itemize}
    \item 数据录入错误
    \item 测量误差
    \item 数据传输错误
    \item 真实的极端情况（较少）
\end{itemize}
\end{knowledgebox}

\subsection{检测异常值}

\begin{lstlisting}[caption={使用IQR方法检测异常值}]
# 计算IQR
Q1 = df_order['quantity'].quantile(0.25)
Q3 = df_order['quantity'].quantile(0.75)
IQR = Q3 - Q1

# 定义异常值范围
lower_bound = Q1 - 1.5 * IQR
upper_bound = Q3 + 1.5 * IQR

# 找出异常值
outliers = df_order[(df_order['quantity'] < lower_bound) | 
                    (df_order['quantity'] > upper_bound)]
print(f"异常值数量：{len(outliers)}")
\end{lstlisting}

\begin{knowledgebox}[知识点：IQR方法原理]
IQR（Interquartile Range，四分位距）是一种稳健的异常值检测方法：

\begin{enumerate}
    \item 把数据从小到大排序
    \item 找到第25\%位置的数（Q1）和第75\%位置的数（Q3）
    \item 计算IQR = Q3 - Q1
    \item 正常数据范围：[Q1 - 1.5×IQR, Q3 + 1.5×IQR]
    \item 超出这个范围的就是异常值
\end{enumerate}

\textbf{优点}：不受极端值影响，适合偏态分布的数据
\end{knowledgebox}

\subsection{处理异常值}

\begin{lstlisting}[caption={处理异常值}]
# 将负数数量修正为1
df_order.loc[df_order['quantity'] < 0, 'quantity'] = 1
\end{lstlisting}

\begin{knowledgebox}[知识点：处理异常值的策略]
\begin{itemize}
    \item \textbf{删除}：直接删除异常值所在的行
    \item \textbf{修正}：根据业务逻辑修正（如负数改为正数）
    \item \textbf{替换}：用平均值、中位数等替换
    \item \textbf{标记}：不删除，但添加标记列
\end{itemize}
\end{knowledgebox}

%========================================
\section{数据格式转换}
%========================================

\subsection{字符串处理}

\begin{lstlisting}[caption={处理价格列}]
# 查看价格列的格式
print(df_order['price'].head())

# 去除"元"字，转换为数值
df_order['price_num'] = df_order['price'].str.replace('元', '')
df_order['price_num'] = df_order['price_num'].astype(float)
\end{lstlisting}

\begin{knowledgebox}[知识点：str访问器]
在Pandas中，要对字符串列进行操作，需要使用\texttt{.str}访问器：

\begin{itemize}
    \item \texttt{df['列名'].str.方法()}
    \item 不能直接写\texttt{df['列名'].strip()}
\end{itemize}

\textbf{常用字符串方法：}
\begin{itemize}
    \item \texttt{str.replace('旧', '新')}：替换字符
    \item \texttt{str.strip()}：去除前后空格
    \item \texttt{str.lower()/upper()}：转换大小写
    \item \texttt{str.contains('关键词')}：判断是否包含
\end{itemize}
\end{knowledgebox}

\subsection{处理手机号}

\begin{lstlisting}[caption={标准化手机号}]
# 去除横线
df_order['phone_clean'] = df_order['phone'].str.replace('-', '')
\end{lstlisting}

\begin{practicebox}[动手练习6：清洗格式混乱数据]
\textbf{任务：}使用04\_数据格式清洗\_练习.csv，清洗姓名和收入列。

\begin{lstlisting}
import pandas as pd

df = pd.read_csv('04_数据格式清洗_练习.csv')

# 查看原始数据
print("姓名处理前：")
print(df['name'].head())

# 去除姓名中的空格
df['name_clean'] = df['name'].str.strip()

print("\n姓名处理后：")
print(df['name_clean'].head())

# 处理收入（将12k转换为12000）
def convert_income(income_str):
    income_str = str(income_str).strip()
    if 'k' in income_str.lower():
        number = float(income_str.lower().replace('k', ''))
        return number * 1000
    else:
        return float(income_str)

df['income_num'] = df['income'].apply(convert_income)

print("\n收入转换：")
print(df[['income', 'income_num']].head(10))
\end{lstlisting}
\end{practicebox}

%========================================
\section{数据分组与统计}
%========================================

\subsection{什么是分组统计？}

\begin{knowledgebox}[分组统计介绍]
分组统计就是：
\begin{enumerate}
    \item 先把数据按某个列分成若干组
    \item 然后对每组进行统计计算
\end{enumerate}

\textbf{类比理解：}
\begin{itemize}
    \item 把全班同学按\textbf{性别}分组，计算每组平均成绩
    \item 把商品按\textbf{类别}分组，计算每类总销量
\end{itemize}
\end{knowledgebox}

\subsection{分组统计}

\begin{lstlisting}[caption={按部门统计工资}]
# 按部门分组，计算平均工资
dept_avg = df_missing.groupby('department')['salary'].mean()
print(dept_avg)

# 多列分组统计
stats = df_missing.groupby('department').agg({
    'salary': ['mean', 'max', 'min'],
    'age': 'mean'
})
print(stats)
\end{lstlisting}

\begin{knowledgebox}[知识点：groupby的工作流程]
\texttt{groupby()}的工作流程：

\begin{enumerate}
    \item \textbf{分组}：按指定列的值把数据分成多组
    \item \textbf{应用}：对每组应用统计函数
    \item \textbf{合并}：把各组结果合并成一个表格
\end{enumerate}

\textbf{常用统计函数：}
\begin{itemize}
    \item \texttt{mean()}：平均值
    \item \texttt{sum()}：求和
    \item \texttt{count()}：计数
    \item \texttt{max()/min()}：最大/最小值
    \item \texttt{std()}：标准差
\end{itemize}
\end{knowledgebox}

\begin{practicebox}[动手练习7：统计电商数据]
\textbf{任务：}使用01\_电商订单数据\_练习.csv，统计各商品的总销量。

\begin{lstlisting}
import pandas as pd

df = pd.read_csv('01_电商订单数据_练习.csv')

# 处理价格（去除"元"字）
df['price_num'] = df['price'].str.replace('元', '').astype(float)

# 计算订单金额
df['total_amount'] = df['price_num'] * df['quantity']

# 按商品分组统计
product_stats = df.groupby('product_name').agg({
    'quantity': 'sum',
    'total_amount': 'sum'
}).sort_values('total_amount', ascending=False)

print("商品销售统计TOP10：")
print(product_stats.head(10))
\end{lstlisting}
\end{practicebox}

%========================================
\section{综合实战：完整数据处理流程}
%========================================

\subsection{案例：完整的电商订单数据处理}

\begin{lstlisting}[caption={完整的数据处理流程}]
import pandas as pd
import numpy as np

# 第1步：读取数据
df = pd.read_csv('01_电商订单数据_练习.csv')
print(f"原始数据：{len(df)}行")

# 第2步：删除重复值
df = df.drop_duplicates()
print(f"删除重复后：{len(df)}行")

# 第3步：处理缺失值
df['user_id'] = df['user_id'].fillna('UNKNOWN')

# 第4步：处理异常值（负数数量）
df.loc[df['quantity'] < 0, 'quantity'] = 1

# 第5步：格式转换（价格）
df['price_num'] = df['price'].str.replace('元', '').astype(float)

# 第6步：格式转换（手机号）
df['phone_clean'] = df['phone'].str.replace('-', '')

# 第7步：计算订单金额
df['total_amount'] = df['price_num'] * df['quantity']

# 第8步：统计分析
product_sales = df.groupby('product_name')['total_amount'].sum()
print("\n各商品销售额：")
print(product_sales.sort_values(ascending=False).head())

# 第9步：保存结果
df.to_csv('清洗后的订单数据.csv', index=False)
print("\n数据已保存到：清洗后的订单数据.csv")
\end{lstlisting}

%========================================
\section{竞赛模拟题}
%========================================

\subsection{题目一：电商订单数据清洗}

\begin{practicebox}[竞赛模拟题1]
\textbf{题目要求：}

使用\textbf{01\_电商订单数据\_练习.csv}，完成以下任务：

\begin{enumerate}
    \item 读取数据，查看基本信息（形状、列名、前5行）
    \item 统计缺失值数量
    \item 删除重复订单
    \item 处理user\_id缺失值（用'UNKNOWN'填充）
    \item 处理负数数量（修正为1）
    \item 提取价格中的数字，转换为数值型
    \item 标准化手机号格式（去除横线）
    \item 计算每个订单的总金额（价格×数量）
    \item 统计每个商品的总销售额和总销量
    \item 保存清洗后的数据到新的CSV文件
\end{enumerate}

\textbf{参考答案：}
\begin{lstlisting}
import pandas as pd

# 读取数据
df = pd.read_csv('01_电商订单数据_练习.csv')

# 查看基本信息
print(f"数据形状：{df.shape}")
print(f"列名：{df.columns.tolist()}")
print(df.head())

# 统计缺失值
print("\n缺失值统计：")
print(df.isnull().sum())

# 删除重复值
df = df.drop_duplicates()

# 填充缺失值
df['user_id'] = df['user_id'].fillna('UNKNOWN')

# 处理异常值
df.loc[df['quantity'] < 0, 'quantity'] = 1

# 格式转换
df['price_num'] = df['price'].str.replace('元', '').astype(float)
df['phone_clean'] = df['phone'].str.replace('-', '')

# 计算金额
df['total_amount'] = df['price_num'] * df['quantity']

# 统计分析
product_stats = df.groupby('product_name').agg({
    'total_amount': 'sum',
    'quantity': 'sum'
}).sort_values('total_amount', ascending=False)

print("\n商品销售统计：")
print(product_stats.head(10))

# 保存结果
df.to_csv('清洗后的订单数据.csv', index=False)
\end{lstlisting}
\end{practicebox}

\subsection{题目二：员工信息数据处理}

\begin{practicebox}[竞赛模拟题2]
\textbf{题目要求：}

使用\textbf{03\_缺失值处理\_练习.csv}，完成以下任务：

\begin{enumerate}
    \item 读取数据，查看缺失值情况
    \item 用中位数填充age列的缺失值
    \item 用平均值填充salary列的缺失值
    \item 用众数填充department列的缺失值
    \item 用众数填充gender列的缺失值
    \item 删除score列仍有缺失值的行
    \item 按department分组，统计平均工资和平均年龄
    \item 保存处理后的数据
\end{enumerate}

\textbf{参考答案：}
\begin{lstlisting}
import pandas as pd

# 读取数据
df = pd.read_csv('03_缺失值处理_练习.csv')

# 查看缺失值
print("缺失值统计：")
print(df.isnull().sum())

# 用中位数填充年龄
median_age = df['age'].median()
df['age'] = df['age'].fillna(median_age)

# 用平均值填充工资
mean_salary = df['salary'].mean()
df['salary'] = df['salary'].fillna(mean_salary)

# 用众数填充部门和性别
mode_dept = df['department'].mode()[0]
df['department'] = df['department'].fillna(mode_dept)

mode_gender = df['gender'].mode()[0]
df['gender'] = df['gender'].fillna(mode_gender)

# 删除score缺失的行
df = df.dropna(subset=['score'])

# 按部门统计
dept_stats = df.groupby('department').agg({
    'salary': 'mean',
    'age': 'mean'
}).round(2)

print("\n部门统计：")
print(dept_stats)

# 保存结果
df.to_csv('处理后的员工数据.csv', index=False)
\end{lstlisting}
\end{practicebox}

%========================================
\section{代码速查表}
%========================================

\subsection{文件操作}

\begin{lstlisting}[caption={读取和保存CSV}]
# 读取CSV
df = pd.read_csv('文件名.csv')

# 保存CSV
df.to_csv('新文件名.csv', index=False)
\end{lstlisting}

\subsection{数据查看}

\begin{lstlisting}[caption={数据查看}]
df.shape          # 形状（行数，列数）
df.columns        # 列名
df.head()         # 前5行
df.tail()         # 后5行
df.info()         # 详细信息
df.describe()     # 统计摘要
df.dtypes         # 数据类型
\end{lstlisting}

\subsection{数据选择}

\begin{lstlisting}[caption={数据选择}]
# 选择列
df['列名']
df[['列1', '列2']]

# 选择行
df.iloc[0]        # 第1行
df.iloc[0:5]      # 前5行

# 条件筛选
df[df['列名'] > 值]
\end{lstlisting}

\subsection{数据清洗}

\begin{lstlisting}[caption={数据清洗}]
# 缺失值
df.isnull().sum()                 # 统计缺失值
df['列'].fillna(值)              # 填充缺失值
df.dropna()                       # 删除缺失值

# 重复值
df.duplicated().sum()             # 统计重复值
df.drop_duplicates()              # 删除重复值

# 异常值（IQR方法）
Q1 = df['列'].quantile(0.25)
Q3 = df['列'].quantile(0.75)
IQR = Q3 - Q1
outliers = df[(df['列'] < Q1-1.5*IQR) | (df['列'] > Q3+1.5*IQR)]
\end{lstlisting}

\subsection{数据转换}

\begin{lstlisting}[caption={数据转换}]
# 类型转换
df['列'].astype(float)
df['列'].astype(str)

# 字符串处理
df['列'].str.replace('旧', '新')
df['列'].str.strip()

# 应用函数
df['新列'] = df['列'].apply(函数)
\end{lstlisting}

\subsection{分组统计}

\begin{lstlisting}[caption={分组统计}]
# 简单分组
df.groupby('分组列')['统计列'].mean()

# 多统计量
df.groupby('分组列').agg({
    '列1': 'mean',
    '列2': ['mean', 'sum']
})
\end{lstlisting}

\end{document}
